% =====================================================================
%  Промежуточный отчёт. Этап 2 ВКР
%
%  Тема: Методика обнаружения несанкционированных действий и аномальной
%        активности в корпоративной сети с использованием нейронных сетей
%
%  Сборка: xelatex -> biber -> xelatex -> xelatex
%  Преамбула и библиография переиспользуются из ../Stage1
% =====================================================================
\documentclass[14pt, a4paper]{extarticle}

% =====================================================================
%  Подключаемые пакеты LaTeX
%  Сборка: XeLaTeX + Biber
%  Соответствие: ГОСТ 7.32-2017, ГОСТ Р 7.0.100-2018
% =====================================================================

% Кодировка и язык
\usepackage{fontspec}
\setmainfont{Times New Roman}
\setsansfont{Arial}
\setmonofont{Courier New}

\usepackage{polyglossia}
\setdefaultlanguage{russian}
\setotherlanguage{english}

% Геометрия страницы (ГОСТ 7.32-2017, п. 6.1):
% левое 30 мм, правое 15 мм, верхнее 20 мм, нижнее 20 мм
\usepackage[
    a4paper,
    left=30mm,
    right=15mm,
    top=20mm,
    bottom=20mm,
    footskip=10mm
]{geometry}

% Межстрочный интервал 1.5 (ГОСТ 7.32-2017, п. 6.3)
\usepackage{setspace}
\onehalfspacing

% Абзацный отступ 1.25 см
\usepackage{indentfirst}
\setlength{\parindent}{1.25cm}

% Математика
\usepackage{amsmath}
\usepackage{amssymb}
\usepackage{amsthm}
\usepackage{mathtools}

% Графика и таблицы
\usepackage{graphicx}
\usepackage{booktabs}
\usepackage{longtable}
\usepackage{tabularx}
\usepackage{array}
\usepackage{multirow}
\usepackage{makecell}
\usepackage{caption}
\usepackage{subcaption}
\usepackage{float}

% Цвет, ссылки, листинги
\usepackage{xcolor}
\usepackage[hidelinks]{hyperref}
\hypersetup{
    pdfencoding=auto,
    bookmarksnumbered=true,
    bookmarksopen=true,
    colorlinks=false
}
\usepackage{listings}
\lstset{
    basicstyle=\ttfamily\footnotesize,
    breaklines=true,
    showstringspaces=false,
    keywordstyle=\bfseries,
    columns=fullflexible,
    captionpos=b,
    frame=single,
    aboveskip=8pt,
    belowskip=8pt
}

% Перечисления и списки
\usepackage{enumitem}
\setlist{nosep, leftmargin=*}

% Заголовки и оглавление
\usepackage{titlesec}
\usepackage{tocloft}

% Библиография по ГОСТ Р 7.0.100-2018
\usepackage[
    backend=biber,
    style=gost-numeric,
    language=auto,
    sorting=none,
    maxbibnames=4,
    minbibnames=1
]{biblatex}

% Дополнительные служебные пакеты
\usepackage{csquotes}
\usepackage{etoolbox}
\usepackage{url}
\usepackage{microtype}
\usepackage{xurl}

% =====================================================================
%  Оформление по ГОСТ 7.32-2017
%  Структурные элементы: разделы, рисунки, таблицы, формулы, списки
% =====================================================================

% Размер основного шрифта 14 pt с межстрочным интервалом 1.5 устанавливается
% в classfile (extarticle/extreport, базовый размер 14pt). Здесь — стилевые
% настройки заголовков и нумерации.

% Заголовки по ГОСТ 7.32-2017 (п. 6.2):
%   - заголовки разделов и подразделов с абзацного отступа,
%     с прописной буквы, без точки в конце, не подчёркиваются;
%   - расстояние между заголовком и текстом — не менее одного интервала;
%   - заголовки разделов набираются жирным шрифтом, размер 14 pt.

\titleformat{\section}
    {\normalfont\fontsize{14}{17}\bfseries}
    {\thesection}{1em}{}
\titlespacing*{\section}{\parindent}{18pt}{12pt}

\titleformat{\subsection}
    {\normalfont\fontsize{14}{17}\bfseries}
    {\thesubsection}{1em}{}
\titlespacing*{\subsection}{\parindent}{14pt}{10pt}

\titleformat{\subsubsection}
    {\normalfont\fontsize{14}{17}\bfseries}
    {\thesubsubsection}{1em}{}
\titlespacing*{\subsubsection}{\parindent}{12pt}{8pt}

% Нумерация рисунков и таблиц — сквозная в пределах раздела
\renewcommand{\thefigure}{\arabic{section}.\arabic{figure}}
\renewcommand{\thetable}{\arabic{section}.\arabic{table}}
\renewcommand{\theequation}{\arabic{section}.\arabic{equation}}

% Подписи рисунков и таблиц по ГОСТ:
%   - «Рисунок N — Название» (тире), снизу, по центру;
%   - «Таблица N — Название», сверху, слева.
\captionsetup[figure]{
    labelsep=endash,
    name=Рисунок,
    justification=centering,
    singlelinecheck=true,
    font=normalsize
}
\captionsetup[table]{
    labelsep=endash,
    name=Таблица,
    justification=raggedright,
    singlelinecheck=false,
    font=normalsize
}

% Колонтитулы — простые, номер страницы по центру нижнего поля
\usepackage{fancyhdr}
\pagestyle{fancy}
\fancyhf{}
\fancyfoot[C]{\thepage}
\renewcommand{\headrulewidth}{0pt}
\renewcommand{\footrulewidth}{0pt}

% Оглавление: «СОДЕРЖАНИЕ» по центру, прописными
\renewcommand{\contentsname}{\hfill СОДЕРЖАНИЕ\hfill}

% Список литературы: «СПИСОК ИСПОЛЬЗОВАННЫХ ИСТОЧНИКОВ»
\defbibheading{bibliography}[Список использованных источников]{%
    \section*{#1}\addcontentsline{toc}{section}{#1}%
}

% =====================================================================
%  Пользовательские команды и сокращения
% =====================================================================

% Краткие ссылки
\newcommand{\figref}[1]{рисунок~\ref{#1}}
\newcommand{\Figref}[1]{Рисунок~\ref{#1}}
\newcommand{\tabref}[1]{таблица~\ref{#1}}
\newcommand{\Tabref}[1]{Таблица~\ref{#1}}
\newcommand{\eqnref}[1]{(\ref{#1})}
\newcommand{\secref}[1]{раздел~\ref{#1}}

% Английские сокращения как термины
\newcommand{\eng}[1]{\textit{#1}}

% Метки общих обозначений
\DeclareMathOperator*{\argmax}{arg\,max}
\DeclareMathOperator*{\argmin}{arg\,min}
\DeclareMathOperator{\softmax}{softmax}
\DeclareMathOperator{\sigm}{\sigma}
\DeclareMathOperator{\KL}{KL}
\DeclareMathOperator{\MMD}{MMD}
\DeclareMathOperator{\PSI}{PSI}

% Векторы и матрицы — обычным начертанием в соответствии с ГОСТ 7.32
% (без полужирного, как принято в большинстве отечественных НИР)

% Окружение для алгоритмов
\newtheoremstyle{algodef}{8pt}{8pt}{\normalfont}{}{\bfseries}{.}{0.5em}{}
\theoremstyle{algodef}
\newtheorem{algorithm}{Алгоритм}[section]
\newtheorem{definition}{Определение}[section]
\newtheorem{statement}{Утверждение}[section]


\addbibresource{../Stage1/bib/refs.bib}

\hypersetup{
    pdftitle={Промежуточный отчёт. Этап 2 ВКР},
    pdfauthor={},
    pdfsubject={Теоретическое сравнение архитектур и проектирование методики},
    pdfkeywords={CNN-LSTM, NIDS, UNSW-NB15, CICIDS2017,
                 ИИ-злоумышленник, dirft, протокол экспериментов}
}

\begin{document}

% ----------- Титульный лист -----------
\begin{titlepage}
\centering
\vspace*{1cm}

{\large\bfseries Промежуточный отчёт по выпускной квалификационной работе}\\[0.3cm]
{\large\bfseries Этап 2. Теоретическое сравнение архитектур и проектирование методики}\\[2cm]

{\Large\bfseries Тема: <<Методика обнаружения несанкционированных действий
и аномальной активности в корпоративной сети с использованием нейронных
сетей>>}\\[3cm]

\begin{flushright}
\begin{tabular}{rl}
Исполнитель: & \rule{6cm}{0.4pt} \\[0.3cm]
Научный руководитель: & \rule{6cm}{0.4pt} \\
\end{tabular}
\end{flushright}

\vfill
{\large \the\year}\\
\end{titlepage}

% ----------- Реферат -----------
\section*{РЕФЕРАТ}
\addcontentsline{toc}{section}{РЕФЕРАТ}

\textbf{Ключевые слова:} нейронные сети, CNN-LSTM, BiLSTM, Transformer-encoder,
автоэнкодер, NIDS, UNSW-NB15, CICIDS2017, sliding windows, focal loss,
drift, PSI, MMD, протокол экспериментов, ИИ-злоумышленник.

\textbf{Объект исследования} --- архитектуры нейронных сетей для
обнаружения аномалий в сетевом трафике и связанные с ними подсистемы
обработки данных и активного тестирования.

\textbf{Предмет исследования} --- проектные решения по выбору архитектуры,
формированию данных и протокола экспериментов для разрабатываемой
методики обнаружения несанкционированных действий.

\textbf{Цель этапа} --- по результатам теоретического анализа и
систематизации опубликованных результатов обосновать выбор архитектуры
нейронной сети, формализовать модель данных и pipeline обработки и
зафиксировать протокол экспериментов Этапа~4 ВКР.

\textbf{Содержание этапа.} Выполнены: систематический анализ
опубликованных результатов применения восьми типов архитектур нейронных
сетей в NIDS-задачах, сводное сопоставление с указанием первоисточников;
формализация теоретической модели данных (схема UNSW-NB15, общие
признаки с CICIDS2017, препроцессинг, окна, балансировка); формирование
библиотеки синтетических атакующих сценариев и протокола экспериментов
E1--E6 с критериями принятия; обоснован выбор CNN-LSTM как основной
модели и зафиксированы альтернативы и baselines.

\newpage

% ----------- Перечень сокращений -----------
\section*{ПЕРЕЧЕНЬ СОКРАЩЕНИЙ И УСЛОВНЫХ ОБОЗНАЧЕНИЙ}
\addcontentsline{toc}{section}{ПЕРЕЧЕНЬ СОКРАЩЕНИЙ И УСЛОВНЫХ ОБОЗНАЧЕНИЙ}

Используются те же сокращения, что и в отчёте Этапа~1; дополнительно
введены:
\begin{description}[font=\normalfont\bfseries, leftmargin=2.5em]
    \item[CI] Confidence Interval (доверительный интервал)
    \item[EDA] Exploratory Data Analysis
    \item[FAR] False Alarm Rate
    \item[FSM] Finite-State Machine (конечный автомат)
    \item[IAT] Inter-Arrival Time (межпакетный интервал)
    \item[IQR] Interquartile Range (межквартильный размах)
    \item[KS-test] Kolmogorov--Smirnov test
    \item[NPZ] NumPy compressed archive
    \item[TSTR] Train Synthetic, Test Real
\end{description}

\newpage

% ----------- Содержание -----------
\tableofcontents
\newpage

% ----------- Введение -----------
\section*{ВВЕДЕНИЕ}
\addcontentsline{toc}{section}{ВВЕДЕНИЕ}

Данный отчёт является продолжением Этапа~1 ВКР и содержит результаты
теоретического сравнения архитектур, проектирования методики обработки
данных и формирования протокола экспериментов. В соответствии с
постановкой задачи (\textit{ВКР.txt}, Этап~2) и проектным решением,
зафиксированным в Этапе~1 (\textit{$H_0^{\text{main}}$}, NFR-1\dots NFR-9),
в качестве основной (\emph{proposed}) архитектуры рассматривается
CNN-LSTM, основные альтернативы --- BiLSTM и Transformer-encoder,
контрольные semi-supervised baselines --- AE и VAE, классические baselines
--- LR/RF/XGBoost/IF/OCSVM.

Особенностью отчёта является \textbf{явная привязка количественных
результатов к первоисточникам}: в таблицах сравнения архитектур
приводятся опубликованные значения метрик с указанием конкретных
рецензируемых работ, использованных датасетов и постановок задачи.
Прямое количественное сравнение между источниками не производится в силу
зафиксированных методологических ограничений; такое сравнение
выполняется в собственном эксперименте Этапа~4 в равных условиях.

\textbf{Структура отчёта.} Раздел~1 содержит сравнительный анализ
архитектур нейронных сетей и обоснование выбора CNN-LSTM с явной
привязкой к опубликованным NIDS-результатам. Раздел~2 формализует
теоретическую модель данных и pipeline обработки. Раздел~3 описывает
библиотеку синтетических сценариев, протокол валидации генератора и
протокол экспериментов E1--E6.

\newpage

% ----------- Основная часть -----------
% =====================================================================
%  Stage 2, раздел 1.
%  Сравнительный анализ архитектур нейронных сетей для NIDS.
%  Опубликованные результаты применения архитектур
%  с явным указанием первоисточников.
% =====================================================================

\section{Сравнительный анализ архитектур нейронных сетей для обнаружения вторжений}

\subsection{Принципы и методология сравнения}

Цель раздела --- сопоставить архитектуры нейронных сетей, рассмотренные в
теоретической части (Этап~1, раздел~1.4), на основе экспериментальных
результатов, опубликованных в рецензируемых источниках в применениях,
структурно эквивалентных задаче настоящей работы --- бинарной/мультиклассовой
детекции аномалий в flow-представлении сетевого трафика.

Для каждой архитектуры приводятся:
\begin{itemize}
    \item первоисточник (научная работа, в которой опубликован результат);
    \item использованный датасет;
    \item постановка задачи (\emph{binary} либо \emph{multiclass});
    \item конкретное количественное значение метрики, опубликованное в работе.
\end{itemize}

\textbf{Методологическое замечание.} Прямое количественное сравнение
результатов из разных источников некорректно по следующим причинам.
\begin{enumerate}
    \item Используются разные датасеты (NSL-KDD, KDD Cup 99, UNSW-NB15,
    CICIDS2017) с принципиально различающимися моделями угроз и распределениями
    признаков~\cite{moustafa2016eval, sommer2010outside, sharafaldin2018toward}.
    \item Используются различные схемы препроцессинга, селекции признаков и
    разделения выборок~\cite{khraisat2019survey, ferrag2020dlcyber}.
    \item Метрики (\emph{Accuracy, Detection Rate, F1, Precision, Recall, AUC})
    интерпретируются в разных задачах различно; в условиях сильного
    дисбаланса классов \emph{Accuracy} даёт оптимистическую
    оценку~\cite{khraisat2019survey, liu2019survey}.
\end{enumerate}

Поэтому опубликованные результаты используются здесь как качественные
индикаторы пригодности архитектуры в NIDS-контексте и для теоретического
обоснования предположения о работоспособности модели на UNSW-NB15. Прямое
количественное сравнение моделей-кандидатов в равных условиях будет
выполнено в собственном эксперименте Этапа~4 настоящей работы.

\subsection{Расширенная теоретическая база сравниваемых архитектур}

В раздел~1 Этапа~1 вошло базовое описание архитектур; здесь приводятся
дополнительные характеристики, существенные для NIDS-применения.

\subsubsection{Полносвязные сети (DNN/MLP)}

В контексте flow-данных DNN рассматривается как нелинейный классификатор
над агрегированными признаками отдельного flow-вектора. Стандартная
композиция включает 2--5 скрытых слоёв с ReLU-активациями и dropout-регуляризацией.
Глубина свыше пяти слоёв в задачах NIDS, по данным
обзоров~\cite{vinayakumar2019dl, ferrag2020dlcyber, ahmad2021nids}, не
приводит к устойчивому улучшению качества.

\subsubsection{1D-CNN над окнами flow-векторов}

Приложение свёрточных слоёв в NIDS интерпретируется как извлечение
\emph{мотивов} (\eng{motifs}) во времени или в признаковой строке.
Стандартная конфигурация: один--три блока \texttt{Conv1D-BatchNorm-ReLU}
с пуллингом, размер ядра \(K \in \{3, 5, 7\}\). При обработке
последовательности окон используется ширина окна
\(W \in \{8, 16, 32, 64\}\), что согласуется с эксплуатационными
характеристиками IPFIX-коллекторов~\cite{rfc7011}.

\subsubsection{Рекуррентные сети: LSTM, GRU, BiLSTM}

В flow-представлении один временной шаг соответствует одному соединению или
агрегированному окну. LSTM/GRU моделируют долгосрочные зависимости между
соединениями, что востребовано для атак с нерегулярной временной структурой
(slow brute-force, low-and-slow exfiltration). BiLSTM использует контекст в
обоих направлениях, что даёт прирост качества за счёт повышенной задержки.

\subsubsection{Гибрид CNN--LSTM}

Иерархия \texttt{Conv1D $\to$ Pooling $\to$ LSTM $\to$ FC} извлекает локальные
паттерны (CNN) и моделирует их временную динамику (LSTM). Гибрид является
устоявшейся и активно публикуемой
архитектурой~\cite{vinayakumar2019dl, ferrag2020dlcyber, ahmad2021nids}.

\subsubsection{TCN}

\emph{Temporal Convolutional Network}~\cite{bai2018tcn} реализует
последовательное моделирование на основе причинных дилатированных свёрток.
В NIDS-публикациях TCN встречается реже LSTM/CNN-LSTM; преимущество ---
параллелизуемость; ограничение --- фиксированное рецептивное поле.

\subsubsection{Transformer-encoder}

Применение Transformer к flow-последовательностям требует позиционного
кодирования и достаточного объёма данных для устойчивого обучения. В
NIDS-публикациях Transformer применяется реже CNN-LSTM;
обзоры~\cite{ferrag2020dlcyber, ahmad2021nids} отмечают сопоставимое или
несколько лучшее качество, но повышенную дисперсию по seed.

\subsubsection{Autoencoder и VAE}

Автоэнкодер обучается на нормальном трафике и идентифицирует аномалии
через ошибку реконструкции; такой режим (\emph{semi-supervised}) ценен
тем, что не требует разметки атак. Эталонная online-реализация для
NIDS --- Kitsune~\cite{mirsky2018kitsune} --- использует ансамбль
автоэнкодеров с инкрементальным обучением.

\subsection{Опубликованные результаты применения DNN/MLP}

В \tabref{tab:dnn-results} сведены опубликованные результаты применения
DNN/MLP-моделей к NIDS-датасетам. Источниками являются рецензируемые работы
по обнаружению вторжений на основе глубокого обучения.

\begin{table}[H]
\centering
\caption{Опубликованные результаты применения DNN/MLP в NIDS}
\label{tab:dnn-results}
\renewcommand{\arraystretch}{1.25}
\setlength{\tabcolsep}{4pt}
\small
\begin{tabularx}{\linewidth}{@{}lXlXX@{}}
\toprule
\textbf{Источник} & \textbf{Архитектура} & \textbf{Датасет} &
\textbf{Постановка} & \textbf{Опубликованный результат} \\
\midrule
Vinayakumar и соавт.~\cite{vinayakumar2019dl} &
DNN, 1--5 скрытых слоёв (\eng{DNN1\dots DNN5}) &
NSL-KDD, KDD Cup 99, UNSW-NB15, CICIDS2017, WSN-DS &
binary и multiclass &
В зависимости от датасета и числа слоёв; авторы рекомендуют DNN из 3--5 слоёв
как универсальную конфигурацию по совокупности accuracy, precision, recall \\
\addlinespace
Sharafaldin и соавт.~\cite{sharafaldin2018toward} &
MLP &
CICIDS2017 &
multiclass &
Precision $\approx 0{,}77$; Recall $\approx 0{,}83$
(уступает RF/ID3, опубликованным в той же работе) \\
\addlinespace
Moustafa и Slay~\cite{moustafa2016eval} &
ANN &
UNSW-NB15 &
binary &
Accuracy $\approx 81{,}3\%$; FAR $\approx 18{,}4\%$ \\
\addlinespace
Khraisat и соавт.~\cite{khraisat2019survey} (обзор) &
DNN &
NSL-KDD, UNSW-NB15 &
binary, multiclass &
Сводные диапазоны Accuracy
$80\text{--}95\%$ для современных DL-моделей; результаты сильно
зависят от препроцессинга \\
\bottomrule
\end{tabularx}
\end{table}

Из \tabref{tab:dnn-results} видно, что DNN/MLP --- адекватная нижняя
граница DL-моделей в NIDS, но устойчиво уступает архитектурам, способным
учитывать локальные и временные структуры (см. далее).

\subsection{Опубликованные результаты применения 1D-CNN}

В \tabref{tab:cnn-results} приведены опубликованные результаты применения
свёрточных архитектур.

\begin{table}[H]
\centering
\caption{Опубликованные результаты применения 1D-CNN в NIDS}
\label{tab:cnn-results}
\renewcommand{\arraystretch}{1.25}
\setlength{\tabcolsep}{4pt}
\small
\begin{tabularx}{\linewidth}{@{}lXlXX@{}}
\toprule
\textbf{Источник} & \textbf{Архитектура} & \textbf{Датасет} &
\textbf{Постановка} & \textbf{Опубликованный результат} \\
\midrule
Vinayakumar и соавт.~\cite{vinayakumar2019dl} &
1D-CNN с пуллингом &
NSL-KDD, UNSW-NB15, KDD Cup 99, CICIDS2017 &
binary и multiclass &
Стабильное превосходство над DNN/MLP по metric F1 на тех же датасетах,
по данным авторов; усиление эффекта при увеличении длины окна \\
\addlinespace
Ferrag и соавт.~\cite{ferrag2020dlcyber} (обзор и собственные эксперименты) &
1D-CNN &
CSE-CIC-IDS2018 и Bot-IoT &
multiclass &
1D-CNN включён в семёрку наиболее перспективных архитектур по сводной
оценке F1 на крупных современных датасетах \\
\addlinespace
Liu и Lang~\cite{liu2019survey} (обзор) &
CNN различных конфигураций &
KDD-семейство, UNSW-NB15 &
binary, multiclass &
Опубликованные значения accuracy в применениях CNN $\geq 90\%$ для большинства
рассмотренных работ; качество выше при использовании CNN-блоков на окнах
последовательных flow \\
\bottomrule
\end{tabularx}
\end{table}

\subsection{Опубликованные результаты применения LSTM/GRU/BiLSTM}

В \tabref{tab:rnn-results} приведены опубликованные результаты для
рекуррентных архитектур.

\begin{table}[H]
\centering
\caption{Опубликованные результаты применения LSTM/GRU/BiLSTM в NIDS}
\label{tab:rnn-results}
\renewcommand{\arraystretch}{1.25}
\setlength{\tabcolsep}{4pt}
\small
\begin{tabularx}{\linewidth}{@{}lXlXX@{}}
\toprule
\textbf{Источник} & \textbf{Архитектура} & \textbf{Датасет} &
\textbf{Постановка} & \textbf{Опубликованный результат} \\
\midrule
Kim и соавт.~\cite{kim2016lstm} &
LSTM-RNN &
KDD Cup 99 &
binary &
Accuracy $\approx 96{,}93\%$; Detection Rate $\approx 98{,}88\%$;
False Alarm Rate $\approx 10{,}04\%$ \\
\addlinespace
Yin и соавт.~\cite{yin2017dlrnn} &
RNN (стандартная) &
NSL-KDD, KDDTest+ &
binary &
Accuracy $\approx 83{,}28\%$ (при 80 hidden units); существенное падение
на KDDTest-21 \\
\addlinespace
Yin и соавт.~\cite{yin2017dlrnn} &
RNN (стандартная) &
NSL-KDD, KDDTest+ &
multiclass &
Accuracy $\approx 81{,}29\%$ \\
\addlinespace
Vinayakumar и соавт.~\cite{vinayakumar2019dl} &
LSTM, GRU &
NSL-KDD, KDD Cup 99, UNSW-NB15 &
binary, multiclass &
LSTM и GRU сопоставимы по средней Accuracy и F1 на NSL-KDD/KDD Cup~99;
GRU стабильно дешевле по числу параметров; на UNSW-NB15 LSTM уступает
гибриду CNN-LSTM в той же работе \\
\addlinespace
Ahmad и соавт.~\cite{ahmad2021nids} (обзор) &
LSTM, BiLSTM &
NSL-KDD, UNSW-NB15, CICIDS2017 &
binary, multiclass &
BiLSTM устойчиво даёт прирост Accuracy/F1 относительно LSTM на NSL-KDD и
UNSW-NB15 за счёт двунаправленного контекста \\
\bottomrule
\end{tabularx}
\end{table}

Особо отметим методологические оговорки:
\begin{itemize}
    \item Высокая Accuracy и DR в~\cite{kim2016lstm} получены на
    KDD Cup 99, корпусе с известными статистическими
    дефектами~\cite{tavallaee2009nslkdd, sommer2010outside}; результат
    рассматривается как историческая точка отсчёта, а не как практически
    переносимая оценка.
    \item Yin и соавт.~\cite{yin2017dlrnn} явно отмечают разрыв качества
    между KDDTest+ и более сложным KDDTest-21, что подчёркивает
    зависимость опубликованных метрик от выборки тестирования.
\end{itemize}

\subsection{Опубликованные результаты применения CNN--LSTM}

В \tabref{tab:cnnlstm-results} сведены опубликованные результаты применения
гибрида CNN--LSTM в задачах NIDS.

\begin{table}[H]
\centering
\caption{Опубликованные результаты применения гибрида CNN--LSTM в NIDS}
\label{tab:cnnlstm-results}
\renewcommand{\arraystretch}{1.25}
\setlength{\tabcolsep}{4pt}
\small
\begin{tabularx}{\linewidth}{@{}lXlXX@{}}
\toprule
\textbf{Источник} & \textbf{Архитектура} & \textbf{Датасет} &
\textbf{Постановка} & \textbf{Опубликованный результат} \\
\midrule
Vinayakumar и соавт.~\cite{vinayakumar2019dl} &
CNN-LSTM &
NSL-KDD, KDD Cup 99, UNSW-NB15, CICIDS2017 &
binary, multiclass &
В семействе DL-моделей данной работы CNN-LSTM показал одни из
наибольших значений F1 и Accuracy на UNSW-NB15 и CICIDS2017
относительно DNN/CNN/LSTM той же работы \\
\addlinespace
Ferrag и соавт.~\cite{ferrag2020dlcyber} &
CNN-LSTM &
CSE-CIC-IDS2018, Bot-IoT, других &
multiclass &
CNN-LSTM фигурирует в перечне эталонных DL-моделей с устойчивыми F1 на
тестовых разделениях; авторы фиксируют преимущество в обнаружении атак с
выраженной временной структурой \\
\addlinespace
Liu и Lang~\cite{liu2019survey} (обзор) &
Гибридные CNN-RNN &
KDD-семейство, UNSW-NB15 &
binary, multiclass &
Гибридные подходы CNN-RNN наиболее устойчиво превосходят как чисто
свёрточные, так и чисто рекуррентные базы по metric F1 в большинстве
рассмотренных работ \\
\addlinespace
Ahmad и соавт.~\cite{ahmad2021nids} (обзор) &
CNN-LSTM, BiLSTM &
NSL-KDD, UNSW-NB15, CICIDS2017 &
binary, multiclass &
В обзоре отмечается, что CNN-LSTM и BiLSTM достигают наилучших Accuracy и F1
в большинстве экспериментов на рассматриваемых датасетах \\
\bottomrule
\end{tabularx}
\end{table}

\subsection{Опубликованные результаты применения автоэнкодеров и VAE}

В \tabref{tab:ae-results} сведены опубликованные результаты применения
автоэнкодерных архитектур в NIDS.

\begin{table}[H]
\centering
\caption{Опубликованные результаты применения AE/VAE в NIDS}
\label{tab:ae-results}
\renewcommand{\arraystretch}{1.25}
\setlength{\tabcolsep}{4pt}
\small
\begin{tabularx}{\linewidth}{@{}lXlXX@{}}
\toprule
\textbf{Источник} & \textbf{Архитектура} & \textbf{Датасет} &
\textbf{Постановка} & \textbf{Опубликованный результат} \\
\midrule
Mirsky и соавт.~\cite{mirsky2018kitsune} &
Kitsune (ансамбль AE с инкрементальным обучением) &
Mirai, OS Scan, SSL Renegotiation, SSDP Flood, ARP MitM, Active Wiretap, Fuzzing, Video Injection (собственный стенд) &
\eng{semi-supervised, online} &
По 9~атакам опубликованы значения AUC, среди которых:
Mirai $\approx 0{,}9988$, SSDP Flood $\approx 0{,}9844$, ARP MitM $\approx 0{,}9606$ \\
\addlinespace
Lopez-Martin и соавт.~\cite{lopez2017cvae} &
Conditional VAE &
NSL-KDD &
multiclass (4 класса) &
По данным авторов CVAE превосходит ряд классических классификаторов
(RF, MLP) на схеме оценки той же работы; конкретные значения см. в первоисточнике \\
\addlinespace
Andresini и соавт.~\cite{andresini2021autoencoder} &
AE-based deep metric learning &
NSL-KDD, AAGM, CICIDS2017 &
binary &
Авторы сообщают о превосходстве предложенного подхода над baseline-AE и
RF/SVM-классификаторами по F1 на трёх датасетах \\
\bottomrule
\end{tabularx}
\end{table}

Замечания по применимости AE/VAE.
\begin{itemize}
    \item Kitsune и его варианты эффективны в \emph{semi-supervised}-режиме
    (обучение только на нормальном трафике), что выгодно отличает их от
    supervised DL-моделей в условиях ограниченной разметки.
    \item Сравнение с supervised-моделями некорректно по принципу: AE
    решают задачу детекции аномалий, в то время как CNN-LSTM и аналоги в
    данной работе нацелены на бинарную/мультиклассовую классификацию с
    разметкой. В настоящей работе AE и VAE используются как контрольные
    semi-supervised baselines.
\end{itemize}

\subsection{Сводное сравнение архитектур по опубликованным результатам}

В \tabref{tab:summary-architectures} приведено сводное сопоставление
архитектур по совокупности качественных характеристик и опубликованных
данных.

\begin{table}[H]
\centering
\caption{Сводное сопоставление архитектур по опубликованным
NIDS-результатам и теоретическим характеристикам}
\label{tab:summary-architectures}
\renewcommand{\arraystretch}{1.25}
\setlength{\tabcolsep}{3pt}
\footnotesize
\begin{tabularx}{\linewidth}{@{}lXccXl@{}}
\toprule
\textbf{Архитектура} & \textbf{Локальные / временные структуры} &
\textbf{Стоимость обучения} & \textbf{Inference latency} &
\textbf{Опубликованные ориентиры качества (NIDS)} & \textbf{Источник} \\
\midrule
DNN/MLP & нет/нет        & низкая   & низкая   &
наибольшая Accuracy на NSL-KDD/UNSW-NB15 в семействе DNN-1\dots5 &
\cite{vinayakumar2019dl, sharafaldin2018toward, moustafa2016eval} \\
\addlinespace
1D-CNN  & да/огр.        & средняя  & низкая   &
устойчивое преимущество над DNN на тех же датасетах в одной и той же работе &
\cite{vinayakumar2019dl, ferrag2020dlcyber, liu2019survey} \\
\addlinespace
LSTM    & нет/да         & средняя  & средняя  &
\(\approx 96{,}93\%\) Acc на KDD Cup 99 в~\cite{kim2016lstm}; уступает
CNN-LSTM на UNSW-NB15 в~\cite{vinayakumar2019dl} &
\cite{kim2016lstm, yin2017dlrnn, vinayakumar2019dl} \\
\addlinespace
GRU     & нет/да         & средняя  & низкая   &
сопоставима с LSTM при меньшем числе параметров &
\cite{vinayakumar2019dl, ahmad2021nids} \\
\addlinespace
BiLSTM  & нет/двунаправ. & высокая  & средняя  &
прирост Accuracy/F1 по сравнению с LSTM на NSL-KDD и UNSW-NB15 &
\cite{ahmad2021nids, ferrag2020dlcyber} \\
\addlinespace
\textbf{CNN-LSTM} & \textbf{да/да} & \textbf{высокая} & \textbf{средняя} &
\textbf{превосходит DNN/CNN/LSTM в одной и той же работе на UNSW-NB15 и CICIDS2017} &
\textbf{\cite{vinayakumar2019dl, ferrag2020dlcyber, liu2019survey, ahmad2021nids}} \\
\addlinespace
TCN         & да/фиксиров. & средняя  & низкая   &
NIDS-публикаций существенно меньше; в обзорах рассматривается как
альтернатива LSTM при ограничениях по latency &
\cite{bai2018tcn, ferrag2020dlcyber} \\
\addlinespace
Transformer & через attention/да & высокая & средняя &
сопоставимое с CNN-LSTM качество в обзорах при повышенной дисперсии по seed &
\cite{vaswani2017attention, ferrag2020dlcyber, ahmad2021nids} \\
\addlinespace
AE          & косвенно/нет  & средняя & низкая &
ансамбль AE Kitsune: AUC $> 0{,}96$ для большинства атак &
\cite{mirsky2018kitsune, andresini2021autoencoder} \\
\addlinespace
VAE         & косвенно/нет  & средняя & низкая &
CVAE превосходит ряд классических baselines на NSL-KDD &
\cite{lopez2017cvae, kingma2014vae} \\
\bottomrule
\end{tabularx}
\end{table}

В \tabref{tab:summary-architectures} обоснованно выделена строка
CNN-LSTM: гибрид сочетает локальное (CNN) и временное (LSTM) моделирование,
а в нескольких независимых рецензируемых
источниках~\cite{vinayakumar2019dl, ferrag2020dlcyber, liu2019survey,
ahmad2021nids} устойчиво демонстрирует одно из лучших значений F1 и
Accuracy в внутренних сравнениях этих работ.

\subsection{Ограничения сравнения и угрозы валидности}

При интерпретации опубликованных результатов следует учитывать следующие
ограничения, систематически отмечаемые в
обзорной литературе~\cite{khraisat2019survey, sommer2010outside,
pendlebury2019tesseract, gharib2016evalframework}.

\begin{enumerate}
    \item \textbf{Различия датасетов.} KDD Cup 99 и NSL-KDD не отражают
    современную модель угроз; результаты на них завышают эффект
    архитектуры и не переносятся на эксплуатационные сценарии. Следовательно,
    высокие метрики LSTM-RNN на KDD Cup 99 (Kim и соавт.,
    \(\approx 96{,}93\%\) Accuracy~\cite{kim2016lstm}) нельзя
    интерпретировать как доказательство практической пригодности.
    \item \textbf{Отсутствие контроля дрейфа.} Большинство работ оценивают
    модели на статически фиксированных тестовых выборках, что приводит к
    \emph{spatial bias}~\cite{pendlebury2019tesseract}; реальная
    устойчивость моделей в потоковом сценарии остаётся неосвещённой.
    \item \textbf{Несбалансированный учёт ложных срабатываний.} Опубликованные
    значения часто относятся к Accuracy, не учитывающей стоимости FP при
    дисбалансе. Работы, приводящие FP-rate, FAR или FP-per-time, остаются
    в меньшинстве.
    \item \textbf{Несравнимость baselines.} Между работами различается состав
    baselines, что затрудняет ранжирование. Например,
    в~\cite{vinayakumar2019dl} и~\cite{ahmad2021nids} CNN-LSTM сравнивается
    с разными конфигурациями DNN/CNN/LSTM, и различие в относительных
    приростах не указывает на устойчивое превосходство.
\end{enumerate}

Эти ограничения определяют необходимость собственного эксперимента в равных
условиях, описание которого приведено в разделе~3 настоящего отчёта (Этап~2)
и реализуется на Этапе~4 ВКР.

\subsection{Обоснование выбора CNN--LSTM в качестве основной модели работы}
\label{subsec:why-cnnlstm}

С учётом теоретического анализа Этапа~1 и систематизированных опубликованных
результатов, представленных в \tabref{tab:summary-architectures}, в качестве
основной (\emph{proposed}) модели работы выбирается CNN--LSTM. Решение
обосновывается следующими аргументами.

\begin{enumerate}
    \item \textbf{Соответствие природе данных.} Атаки в flow-представлении
    проявляются одновременно как локальные паттерны (например, специфическое
    сочетание признаков SYN-флуда в одном flow) и как временные тренды
    (последовательность подобных flow). CNN--LSTM моделирует обе компоненты
    в одной архитектуре~\cite{vinayakumar2019dl, ferrag2020dlcyber}.
    \item \textbf{Опубликованные результаты.} В нескольких независимых
    рецензируемых работах~\cite{vinayakumar2019dl, ferrag2020dlcyber,
    liu2019survey, ahmad2021nids} CNN-LSTM демонстрирует наибольшие F1 и
    Accuracy в \emph{внутреннем} сравнении с DNN, CNN, LSTM и BiLSTM
    на UNSW-NB15 и CICIDS2017.
    \item \textbf{Воспроизводимость.} Архитектура хорошо документирована,
    реализуется стандартными примитивами PyTorch, не требует
    специализированных техник обучения (как Transformer на ограниченных
    выборках), что обеспечивает воспроизводимость экспериментов.
    \item \textbf{Эксплуатационная стоимость.} Inference latency
    CNN-LSTM, согласно обзорам~\cite{ferrag2020dlcyber, ahmad2021nids},
    укладывается в диапазон, приемлемый для NTA-pipeline (десятки
    миллисекунд при batch=1 на CPU); это совместимо с требованием NFR-3
    раздела~4 Этапа~1.
    \item \textbf{Альтернативы и контрольные модели.} В качестве альтернатив
    в собственном эксперименте сохраняются BiLSTM (двунаправленный контекст)
    и Transformer-encoder (attention); в качестве контрольных
    semi-supervised baselines --- AE и VAE; в качестве классических baseline
    --- Logistic Regression, Random Forest, XGBoost, Isolation Forest,
    One-Class SVM.
\end{enumerate}

\subsection{Промежуточные выводы по разделу}

\begin{enumerate}
    \item Систематизированы опубликованные результаты применения восьми
    типов архитектур в NIDS-задачах, эквивалентных постановке настоящей
    работы. Все результаты приведены с указанием источника, датасета и
    постановки.
    \item Зафиксированы три группы методологических ограничений сравнения
    (несовместимость датасетов, отсутствие контроля дрейфа, доминирование
    Accuracy при дисбалансе классов), что определяет необходимость
    собственного эксперимента в равных условиях.
    \item Сводное сопоставление по характеристикам архитектур и опубликованным
    NIDS-результатам подтверждает теоретический выбор CNN-LSTM как основной
    архитектуры; альтернативами в равных условиях останутся BiLSTM и
    Transformer-encoder, контрольными semi-supervised baselines --- AE и VAE,
    классическими baseline --- LR/RF/XGBoost/IF/OCSVM.
    \item Полученное обоснование является \emph{теоретическим} и подтверждается
    \emph{ссылочным} образом; \emph{количественное} подтверждение возможно
    только в собственном эксперименте, выполняемом на Этапе~4 ВКР по
    протоколу, формализованному в разделе~3 настоящего отчёта.
\end{enumerate}

% =====================================================================
%  Stage 2, раздел 2.
%  Теоретическая модель данных и pipeline обработки.
% =====================================================================

\section{Теоретическая модель данных и pipeline обработки}

\subsection{Схема признаков UNSW-NB15}

UNSW-NB15~\cite{moustafa2015unsw, moustafa2016eval, unswportal} в составе
\emph{cleaned}-распределения содержит 49 признаков, разделяемых на пять
групп. Назначение каждой группы и её роль в детекции атак приведены в
\tabref{tab:unsw-feature-groups}.

\begin{table}[H]
\centering
\caption{Группы признаков UNSW-NB15 и их назначение}
\label{tab:unsw-feature-groups}
\renewcommand{\arraystretch}{1.25}
\begin{tabularx}{\linewidth}{@{}lXX@{}}
\toprule
\textbf{Группа} & \textbf{Состав} & \textbf{Роль в детекции} \\
\midrule
Базовые потоковые &
\emph{dur, proto, service, state, sbytes, dbytes, sttl, dttl, sloss, dloss}
&
непосредственное описание flow: длительность, протокол, состояние сессии,
байты и пакеты в обе стороны \\
\addlinespace
Контентные &
\emph{sload, dload, swin, dwin, smean, dmean, stcpb, dtcpb, trans\_depth,
response\_body\_len} &
скорости, окна TCP, базовые транзакционные характеристики \\
\addlinespace
Временные &
\emph{sjit, djit, sintpkt, dintpkt, tcprtt, synack, ackdat} &
межпакетные интервалы, джиттеры, RTT, timing-сигнатуры \\
\addlinespace
Дополнительные сгенерированные &
\emph{ct\_state\_ttl, ct\_flw\_http\_mthd, is\_ftp\_login,
ct\_ftp\_cmd} &
агрегаты по протокольному поведению \\
\addlinespace
Контекстные \emph{ct\_*}&
\emph{ct\_srv\_src, ct\_srv\_dst, ct\_dst\_ltm, ct\_src\_ltm,
ct\_src\_dport\_ltm, ct\_dst\_sport\_ltm, ct\_dst\_src\_ltm} &
агрегаты по последним 100 соединениям; косвенно отражают режим
сканирования и фан-аут \\
\bottomrule
\end{tabularx}
\end{table}

Метки представлены двумя полями: \emph{label} \(\in \{0, 1\}\) (бинарная) и
\emph{attack\_cat} \(\in \mathcal{C} \cup \{\varnothing\}\) (мультиклассовая,
с пустым значением для нормального трафика).

\subsection{Схема признаков CICIDS2017 и подмножество для cross-evaluation}

CICIDS2017~\cite{sharafaldin2018toward} формируется утилитой CICFlowMeter и
содержит более 80 flow-признаков. Часть из них функционально эквивалентна
полям UNSW-NB15 (длительность, протокол, число пакетов и байт в обе
стороны, межпакетные интервалы, скорости передачи). Подмножество общих
признаков, принимаемое в работе, представлено в
\tabref{tab:common-features}.

\begin{table}[H]
\centering
\caption{Подмножество признаков, общих для UNSW-NB15 и CICIDS2017,
принятое для cross-evaluation}
\label{tab:common-features}
\renewcommand{\arraystretch}{1.25}
\begin{tabularx}{\linewidth}{@{}lll@{}}
\toprule
\textbf{Семантика} & \textbf{UNSW-NB15} & \textbf{CICIDS2017 (CICFlowMeter)} \\
\midrule
Длительность потока       & \emph{dur}      & \emph{Flow Duration} \\
Протокол транспорта        & \emph{proto}    & \emph{Protocol} \\
Прикладной сервис          & \emph{service}  & \emph{Destination Port}+\emph{Service heuristics} \\
Байты от источника         & \emph{sbytes}   & \emph{Total Length of Fwd Packets} \\
Байты к источнику          & \emph{dbytes}   & \emph{Total Length of Bwd Packets} \\
Пакеты от источника        & \emph{spkts}    & \emph{Total Fwd Packets} \\
Пакеты к источнику         & \emph{dpkts}    & \emph{Total Backward Packets} \\
TTL источника              & \emph{sttl}     & ---  (косвенно через \emph{Init\_Win\_bytes\_forward})\\
Скорость источника         & \emph{sload}    & \emph{Fwd Packets/s} \\
Скорость получателя        & \emph{dload}    & \emph{Bwd Packets/s} \\
Межпакетный интервал, src  & \emph{sintpkt}  & \emph{Fwd IAT Mean} \\
Межпакетный интервал, dst  & \emph{dintpkt}  & \emph{Bwd IAT Mean} \\
Джиттер, src               & \emph{sjit}     & \emph{Fwd IAT Std} \\
Джиттер, dst               & \emph{djit}     & \emph{Bwd IAT Std} \\
SYN-ACK / ACK-DAT timings  & \emph{synack, ackdat} & \emph{Flow IAT}-агрегаты \\
\bottomrule
\end{tabularx}
\end{table}

Признаки, не имеющие прямого аналога (TTL, ряд \emph{ct\_*}-агрегатов
UNSW-NB15, специфические признаки CICFlowMeter), исключаются из набора
cross-evaluation. Эта позиция фиксирована и применяется в эксперименте E4
протокола (раздел~3 Этапа~2).

\subsection{Pipeline обработки данных}

Полный pipeline обработки данных принимает на вход поток flow-записей и
формирует на выходе скоринг \(s(x)\) и решение \(\hat{y}\):

\[
\boxed{\;
\text{IPFIX/CSV} \to
\text{Schema-Valid} \to
\text{Cleanup} \to
\text{Encoding} \to
\text{Scaling} \to
\text{Windowing} \to
\text{Model} \to
\text{Threshold} \to
\text{Alert}
\;}
\]

Стадия каждого блока формализуется ниже.

\subsection{Очистка данных}

В UNSW-NB15 документированы~\cite{moustafa2015unsw} следующие виды
``шума'': пропуски в полях \emph{ct\_flw\_http\_mthd, is\_ftp\_login,
ct\_ftp\_cmd}, неконсистентные строковые значения \emph{service} и
\emph{state}, отдельные крайние выбросы в \emph{dur, sbytes, dbytes}.
Применяемая очистка:

\begin{enumerate}
    \item Замена нечисловых пропусков константой \texttt{"-"} (отсутствие
    значения). Для числовых пропусков --- замена нулём, что согласуется с
    интерпретацией признаков (отсутствие активности).
    \item Унификация регистра и пробелов в \emph{service}, \emph{state},
    \emph{proto}.
    \item Удаление дубликатов по полному вектору признаков на стороне
    обучающей выборки (исключение перекрёстной утечки).
    \item Идентификация и обрезка выбросов в крайних \(0{,}1\%\) хвостов
    распределений для \emph{dur, sbytes, dbytes, sload, dload}; границы
    обрезки фиксируются по обучающей выборке и применяются к валидационной
    и тестовой без пересчёта.
\end{enumerate}

Замечание: вся статистика для очистки и масштабирования вычисляется
\emph{только} на обучающей выборке, что исключает test-time leakage.

\subsection{Кодирование категориальных признаков}

\subsubsection{One-hot}

Для признака \(c\) с конечным алфавитом \(\mathcal{A}_c =
\{a_1, \dots, a_K\}\) one-hot-вектор имеет вид
\(e(c) = (\mathbb{I}[c=a_1], \dots, \mathbb{I}[c=a_K])^\top \in \{0,1\}^K\).
Применимо для признаков с малой кардинальностью --- \emph{proto, state}.

\subsubsection{Frequency-encoding и target-encoding}

Для признака с большой кардинальностью \emph{service}
(\(\sim 13\) уровней) применяется frequency-encoding:
\begin{equation}
e_{\text{freq}}(c) = \frac{|\{i : x_i^{(c)} = c\}|}{n_{\text{train}}},
\label{eq:freq-enc}
\end{equation}
а как опциональная альтернатива --- target-encoding со сглаживанием:
\begin{equation}
e_{\text{tgt}}(c) =
\frac{\sum_i \mathbb{I}[x_i^{(c)}=c]\,y_i + \alpha\,\bar{y}}
     {\sum_i \mathbb{I}[x_i^{(c)}=c] + \alpha},
\label{eq:tgt-enc}
\end{equation}
где \(\bar{y}\) --- средняя метка класса в обучающей выборке, \(\alpha\) ---
параметр сглаживания. Target-encoding применяется с осторожностью и только
при кросс-валидационной схеме во избежание утечки.

\subsubsection{Embedding categorical}

В DL-моделях категориальные признаки могут кодироваться обучаемыми
embeddings \(E_c \in \mathbb{R}^{K \times d_e}\), где \(d_e\) --- размерность
embedding. Это совместимо с архитектурой CNN-LSTM и, по обзорам
NIDS-DL~\cite{ferrag2020dlcyber, ahmad2021nids}, даёт прирост качества
относительно one-hot для признаков большой кардинальности.

\subsection{Преобразования числовых признаков}

\subsubsection{Логарифмическое преобразование}

Признаки с положительной асимметрией и тяжёлыми хвостами (\emph{dur, sbytes,
dbytes, sload, dload, sintpkt, dintpkt}) подвергаются преобразованию
\begin{equation}
\tilde{x} = \log(1 + x).
\label{eq:log1p}
\end{equation}
\(\log(1+x)\) корректно обрабатывает нулевые значения и эквивалентна
гомеоморфизму, переводящему распределение в близкое к нормальному.

\subsubsection{Robust-scaling}

После логарифмирования признаки масштабируются робастно:
\begin{equation}
\tilde{x}' = \frac{\tilde{x} - \operatorname{med}(\tilde{x})}
                  {\operatorname{IQR}(\tilde{x})},
\label{eq:robust-scaler}
\end{equation}
где \(\operatorname{med}\) --- медиана, \(\operatorname{IQR}\) ---
интерквартильный размах, оцениваемые \emph{на обучающей выборке}.
Робастные оценки устойчивы к выбросам, что критично для NIDS, где
аномалии формируют ``хвосты'' распределения.

\subsubsection{Стандартизация}

Для baseline-моделей (LR, MLP) дополнительно используется стандартизация
\(z = (x - \mu)/\sigma\). Для DL-моделей эта операция эквивалентна robust-scaling
с заменой \(\operatorname{med} \to \mu\), \(\operatorname{IQR} \to \sigma\).

\subsection{Формирование скользящих окон}

Для рекуррентных и гибридных моделей (CNN-LSTM, BiLSTM, Transformer)
исходный поток flow-записей преобразуется в скользящие окна фиксированной
длины \(W\) с шагом \(S\):
\[
\mathcal{W}_t = \bigl(x_{t-W+1}, x_{t-W+2}, \dots, x_t\bigr) \in \mathbb{R}^{W \times d},
\]
где \(d\) --- размерность признакового вектора после препроцессинга и
кодирования.

\subsubsection{Выбор параметров \texorpdfstring{$W$ и $S$}{W и S}}

Параметр \(W\) определяет компромисс между разрешающей способностью по
времени (короткие \(W\) лучше реагируют на мгновенные всплески) и
полнотой контекста (длинные \(W\) лучше учитывают slow-атаки). По данным
обзоров~\cite{vinayakumar2019dl, ferrag2020dlcyber, ahmad2021nids}
и эталонной работы Kim и соавт.~\cite{kim2016lstm} распространённый диапазон
для NIDS-окон в flow-представлении --- \(W \in \{8, 16, 32, 64\}\).
В настоящей работе принят базовый \(W = 32\) с проверкой устойчивости при
\(W \in \{16, 32, 64\}\) в рамках Эксперимента~1 Этапа~4.

Шаг \(S\) для основного режима принят равным \(S = W/4\) (75\% перекрытие).
Это согласуется с практикой потокового анализа в NTA: малое \(S\) уменьшает
TTD, большое --- сокращает вычислительные затраты.

\subsubsection{Целевая разметка окна}

Для бинарной задачи метка окна формируется как
\begin{equation}
y_{\mathcal{W}_t} =
\mathbb{I}\!\left[\sum_{i=t-W+1}^{t} y_i > 0\right],
\label{eq:window-label}
\end{equation}
то есть окно считается атакующим при наличии хотя бы одной атакующей
записи. Этот вариант принят как основной (\emph{any-positive}); как
альтернатива в Эксперименте~1 рассматривается \emph{majority-positive}
(порог по доле атакующих записей).

\subsubsection{Запрет утечки и временной split}

Train/val/test-разделение выполняется при сохранении временного порядка:
обучающая выборка предшествует валидационной, валидационная ---
тестовой. При этом запрещено формирование окна, содержащего записи из
разных сплитов, --- в противном случае возникает утечка
``из будущего''~\cite{pendlebury2019tesseract}.

\subsection{Управление дисбалансом классов}

UNSW-NB15 в \emph{cleaned}-распределении имеет умеренный дисбаланс
классов с доминированием \emph{Generic}. Применяемые меры:

\subsubsection{Re-weighting}

Веса классов \(w_c\) обратно пропорциональны их частотам:
\begin{equation}
w_c = \frac{n}{|\mathcal{C}|\,n_c},
\label{eq:reweight}
\end{equation}
где \(n_c\) --- число примеров класса \(c\) в обучающей выборке. Веса
включаются в функцию потерь.

\subsubsection{Focal loss}

Для усиления внимания к ``тяжёлым'' примерам применяется фокальная функция
потерь~\cite{lin2017focal}:
\begin{equation}
\mathcal{L}_{\text{focal}} =
- \frac{1}{n}\sum_{i=1}^{n}
  \alpha_{y_i}\,(1-p_{y_i})^{\gamma}\,\log p_{y_i},
\label{eq:focal}
\end{equation}
где \(p_{y_i}\) --- предсказанная вероятность истинного класса,
\(\alpha\), \(\gamma\) --- параметры фокусировки. Стандартное значение
\(\gamma = 2{,}0\)~\cite{lin2017focal}.

\subsubsection{SMOTE}

Для классических baseline (LR, RF, XGBoost) применяется
SMOTE~\cite{chawla2002smote}: для каждой точки миноритарного класса
\(x_i\) и её \(k\)-ближайшего соседа того же класса \(x_j\) формируется
синтетическая точка
\begin{equation}
x_{\text{new}} = x_i + \lambda\,(x_j - x_i),\quad \lambda \sim U(0, 1).
\label{eq:smote}
\end{equation}
SMOTE не применяется к окончательной DL-обучающей выборке, так как при
последовательной структуре окон вносит синтетические разрывы во временную
ось.

\subsection{Контроль дрейфа распределений}

В соответствии с разделом~3 Этапа~1 в pipeline предусматривается отдельный
блок drift-monitoring, реализующий метрики:
\begin{itemize}
    \item \(\PSI\) над выбранными признаками (\eqnref{eq:psi});
    \item \(\KL\)-расхождение над дискретизованными распределениями;
    \item \(\MMD\) над непрерывными признаками
    (\eqnref{eq:mmd}).
\end{itemize}

Эталонное окно --- начальное окно нормального трафика в обучающей выборке.
Текущее окно --- скользящее окно того же размера в потоке инференса. При
превышении пороговых значений хотя бы по одной из метрик активируется
расширенный режим логирования и (опционально) переобучение модели.
Конкретные пороги \(\tau_{\PSI} = 0{,}25\), \(\tau_{\MMD}\) (определяется
по 95-перцентилю распределения MMD на bootstrap-выборках обучающего
блока) фиксированы и применяются в Эксперименте~5 Этапа~4.

\subsection{Спецификация форматов входов и выходов}

В \tabref{tab:io-spec} приведены спецификации входных и выходных форматов
системы.

\begin{table}[H]
\centering
\caption{Спецификации входных и выходных данных системы}
\label{tab:io-spec}
\renewcommand{\arraystretch}{1.25}
\begin{tabularx}{\linewidth}{@{}llX@{}}
\toprule
\textbf{Назначение} & \textbf{Формат} & \textbf{Содержание} \\
\midrule
Сырой поток flow-записей &
CSV / JSON Lines &
Поля UNSW-NB15 (49) с типами по \emph{schema.json}; разделитель ``,''; UTF-8 \\
\addlinespace
Препроцессированный поток &
NPZ / Parquet &
Числовая матрица \(\mathbb{R}^{n \times d}\) после кодирования и масштабирования \\
\addlinespace
Окна последовательности &
NPZ / Parquet &
Тензор \(\mathbb{R}^{m \times W \times d}\) с метками \(y \in \{0,1\}^m\) \\
\addlinespace
Скоринг детектора &
JSON Lines &
\texttt{\{ts, score, decision, attack\_cat\_pred, prob, window\_id\}} \\
\addlinespace
Алерт &
JSON Lines / SIEM-event &
\texttt{\{alert\_id, ts, severity, score, profile, sample\_features\}} \\
\addlinespace
Drift-метрика &
JSON Lines &
\texttt{\{ts, psi, kl, mmd, ref\_window, cur\_window\}} \\
\bottomrule
\end{tabularx}
\end{table}

\subsection{Воспроизводимость и контроль конфигурации}

Все шаги pipeline параметризуются конфигурационными YAML-файлами в
директории \texttt{configs/}:

\begin{itemize}
    \item \texttt{configs/data/unsw\_nb15.yaml} --- пути, схема, параметры
    очистки и кодирования;
    \item \texttt{configs/data/cicids2017.yaml} --- то же для CICIDS2017;
    \item \texttt{configs/preprocess/default.yaml} --- параметры
    масштабирования, encoding, windowing;
    \item \texttt{configs/models/cnn\_lstm.yaml} --- архитектура и
    гиперпараметры основной модели;
    \item \texttt{configs/train/default.yaml} --- параметры обучения,
    seed, логи;
    \item \texttt{configs/eval/default.yaml} --- метрики, пороги, схема
    бутстрап-CI.
\end{itemize}

Каждый запуск экспериментов фиксирует:
\begin{enumerate}
    \item версию кода (commit hash);
    \item набор seed (\(\geq 5\));
    \item гиперпараметры обучения и препроцессинга;
    \item версии зависимостей (\texttt{pip freeze});
    \item полное содержимое использованных конфигов в виде attached-артефактов.
\end{enumerate}

\subsection{Промежуточные выводы по разделу}

\begin{enumerate}
    \item Полностью описана схема признаков UNSW-NB15 (5 групп, 49
    признаков) и зафиксировано подмножество, общее с CICIDS2017
    (12--14 признаков), пригодное для протокола cross-evaluation.
    \item Сформирован полный pipeline обработки данных:
    schema-validation $\to$ cleanup $\to$ encoding (one-hot, frequency,
    target, embedding) $\to$ scaling (log, robust) $\to$ windowing
    (\(W=32\), \(S=W/4\)) $\to$ model $\to$ threshold $\to$ alert.
    Каждая стадия снабжена математическим описанием.
    \item Зафиксированы меры управления дисбалансом классов: re-weighting,
    focal loss с параметрами \(\alpha\), \(\gamma\), и SMOTE для классических
    baseline. Введён запрет временной утечки при формировании окон.
    \item Описан блок контроля дрейфа на основе метрик PSI, KL, MMD с
    конкретными пороговыми значениями для эксперимента E5; определены
    эталонное и текущее окна.
    \item Зафиксированы спецификации входных и выходных форматов на всех
    стадиях pipeline и принципы конфигурационного контроля воспроизводимости.
\end{enumerate}

% =====================================================================
%  Stage 2, раздел 3.
%  Синтетические тестовые сценарии и протокол экспериментов.
% =====================================================================

\section{Синтетические тестовые сценарии и протокол экспериментов}

\subsection{Назначение раздела}

В разделе формализованы (а) библиотека синтетических тестовых сценариев,
формируемая ИИ-злоумышленником, и (б) протокол экспериментов Этапа~4 ВКР.
Раздел является связующим звеном между методическими решениями
разделов~1--2 настоящего отчёта (Этап~2) и реализацией прототипа на
Этапе~3.

\subsection{Библиотека синтетических атакующих сценариев}

В соответствии с обоснованием раздела~3 Этапа~1, библиотека реализуется
в виде набора параметризованных шаблонов, согласованных с таксономией
UNSW-NB15. Перечень и параметры шаблонов приведены в
\tabref{tab:attack-templates}.

\begin{table}[H]
\centering
\caption{Библиотека синтетических атакующих шаблонов}
\label{tab:attack-templates}
\renewcommand{\arraystretch}{1.25}
\setlength{\tabcolsep}{4pt}
\small
\begin{tabularx}{\linewidth}{@{}lXX@{}}
\toprule
\textbf{Шаблон} & \textbf{Параметризация} & \textbf{Соответствие классу UNSW-NB15} \\
\midrule
\texttt{ddos\_synflood} &
интенсивность \(\lambda_0\), длительность \(\Delta\), доля SYN-флагов,
распределение исходных IP &
DoS \\
\addlinespace
\texttt{port\_scan} &
скорость подключений \(\rho\), диапазон портов, режим (последовательный /
случайный), длительность сессии &
Reconnaissance \\
\addlinespace
\texttt{brute\_force} &
интервал между попытками, целевой порт (22, 3389, 80, 443),
доля успешных авторизаций &
Reconnaissance / Exploits \\
\addlinespace
\texttt{exploit\_payload} &
размер payload, нестандартные сочетания TCP-флагов, доля retransmissions &
Exploits / Backdoors \\
\addlinespace
\texttt{data\_exfil} &
длительность сессии, асимметрия \(b_{out}/b_{in}\), стабильность скорости,
прикладной сервис &
Backdoors / Worms \\
\addlinespace
\texttt{internal\_recon} &
фан-аут потоков (число destinations за период), периодичность,
размер сессии &
Reconnaissance \\
\addlinespace
\texttt{worm\_lateral} &
число направленных потоков, периодичность, размер сессии,
распределение destinations &
Worms \\
\bottomrule
\end{tabularx}
\end{table}

Каждый шаблон возвращает упорядоченную последовательность flow-векторов,
семантически согласованную с признаковой схемой UNSW-NB15. Параметры по
умолчанию выбраны из обоснованных в литературе диапазонов
характеристик каждого класса
атак~\cite{moustafa2015unsw, sharafaldin2018toward, sharafaldin2019ddos,
khraisat2019survey}.

\subsection{Конечно-автоматная политика агента}

Конечный автомат агента, формализованный в разделе~3 Этапа~1, в текущем
протоколе принимает следующее множество состояний:
\[
\mathcal{S} =
\{\text{NORMAL}\} \cup
\{\text{ATTACK}(c)\}_{c \in \{\,\text{ddos, scan, brute, exploit, exfil, intrec, worm}\,\}} \cup
\{\text{DRIFT}(k)\}_{k \in \{\,\text{cov, prior, concept}\,\}}.
\]

Политика \(\pi\) задаётся вероятностями переходов, минимальными длительностями
состояний и параметрами рандомизации шаблонов.

\subsection{Протокол валидации генератора}

Перед использованием в экспериментах E1--E6 (см. ниже) генератор проходит
валидацию по двум критериям, формализованным в разделе~3 Этапа~1:

\begin{enumerate}
    \item \textbf{Маргинальное соответствие нормы.} Двувыборочный тест
    Колмогорова--Смирнова и MMD-сравнение распределений признаков
    \emph{dur, sbytes, dbytes, sttl, sintpkt, dintpkt} режима NORMAL агента
    с распределениями нормального трафика обучающей выборки UNSW-NB15.
    Условие приёма: $p\text{-value}_{\text{KS}} > 0{,}01$ и MMD не
    превышает 95-перцентиля bootstrap-распределения.
    \item \textbf{Семантическая согласованность атак.} Сравнение средних и
    медиан ключевых признаков (\emph{dur, sbytes, dbytes, sintpkt}) для
    каждого ATTACK-состояния с соответствующим классом UNSW-NB15
    (DoS, Reconnaissance, Exploits, Backdoors, Worms). Условие приёма:
    относительное отклонение средних не более 25\% и совпадение знака
    асимметрии распределений.
\end{enumerate}

При невыполнении условий приёма параметры шаблона корректируются до
выполнения; результаты валидации фиксируются в журнале и приводятся в
отчёте Этапа~4.

\subsection{Протокол экспериментов}

В \tabref{tab:experiments} зафиксированы шесть экспериментов с указанием
цели, входов, метрик и условий принятия гипотезы.

\begin{table}[H]
\centering
\caption{Протокол экспериментов Этапа~4 ВКР}
\label{tab:experiments}
\renewcommand{\arraystretch}{1.25}
\setlength{\tabcolsep}{4pt}
\small
\begin{tabularx}{\linewidth}{@{}lXXX@{}}
\toprule
\textbf{ID} & \textbf{Цель и условия} & \textbf{Метрики} & \textbf{Критерий принятия} \\
\midrule
E1 &
Сравнение моделей-кандидатов (CNN-LSTM, BiLSTM, LSTM, GRU, 1D-CNN, MLP,
Transformer, AE, VAE, RF, XGBoost, LR, IF, OCSVM) на UNSW-NB15;
\(\geq 5\) seed; одинаковый препроцессинг и пороги &
Precision, Recall, F1, MCC, PR-AUC, ROC-AUC; \(\sigma_{F1}\) &
\(F_1^{\text{CNN-LSTM}} \geq 0{,}90\); статистически значимое
($p < 0{,}05$, парный Уилкоксон) превосходство над ансамблем baselines \\
\addlinespace
E2 &
Анализ ошибок (FP/FN) лучшей модели; разрезы по \emph{attack\_cat},
по протоколу, по сервису &
confusion matrix, FN/FP по классам, ROC по классам &
явное отображение классов, в которых F1 ниже среднего значения по тесту \\
\addlinespace
E3 &
Калибровка порога; температурное масштабирование &
ECE, FP-rate vs Recall, FP-per-time, TTD &
наличие рабочей точки, удовлетворяющей NFR-2 \\
\addlinespace
E4 &
Тестирование с ИИ-злоумышленником на стенде; 6 шаблонных сценариев
+ 3 drift-сценария &
Precision, Recall, F1, TTD, FP-per-time;
ground truth из FSM-журнала &
\(\Delta F_1 \leq 0{,}10\) при agent-NORMAL+ATTACK без drift; TTD$\leq W$
для шаблонов с явными flow-проявлениями \\
\addlinespace
E4' &
Cross-evaluation: train UNSW $\to$ test CICIDS2017
без дообучения и c дообучением на 10\% &
Precision, Recall, F1, PR-AUC, $\Delta F_1$ относительно UNSW &
$\Delta F_1 \leq 0{,}20$ без дообучения; $\leq 0{,}10$ с дообучением \\
\addlinespace
E5 &
Drift-устойчивость; PSI/MMD-инжектирование &
PSI, MMD, $\Delta F_1$ vs PSI; время восстановления (recovery time) &
$\Delta F_1 \leq 0{,}10$ при PSI $\leq 0{,}25$; срабатывание
drift-monitor \\
\addlinespace
E6 &
Эксплуатационная производительность &
inference latency (ms / window),
throughput (EPS), CPU/RAM &
NFR-3, NFR-4 \\
\bottomrule
\end{tabularx}
\end{table}

\subsection{Меры обеспечения статистической достоверности}

В соответствии с замечаниями TESSERACT~\cite{pendlebury2019tesseract} и
обзоров~\cite{khraisat2019survey, sommer2010outside} для каждого эксперимента
реализуются:

\begin{enumerate}
    \item \textbf{Множественные seed.} \(n \geq 5\) независимых запусков
    с фиксированными seed, отчётность по среднему и доверительному интервалу
    95\% (bootstrap, \(B = 1000\)).
    \item \textbf{Парные сравнения.} При сравнении пар моделей по seed
    применяется парный критерий Уилкоксона.
    \item \textbf{Запрет утечки.} Train/val/test-разделение по времени;
    препроцессорные статистики --- только по train.
    \item \textbf{Калибровка.} Все модели с поддержкой
    \(\hat{p}\) проходят температурное масштабирование на val-выборке;
    отчётность по ECE до и после.
    \item \textbf{Воспроизводимость.} Версионирование кода (commit hash),
    конфигов и данных; журнал \texttt{run\_id} для каждого запуска.
\end{enumerate}

\subsection{Базовые модели для сравнения}

Базовые модели (\emph{baselines}), используемые в эксперименте E1, и их
роль в работе приведены в \tabref{tab:baselines}.

\begin{table}[H]
\centering
\caption{Базовые модели сравнения, используемые в эксперименте E1}
\label{tab:baselines}
\renewcommand{\arraystretch}{1.25}
\begin{tabularx}{\linewidth}{@{}lXX@{}}
\toprule
\textbf{Модель} & \textbf{Тип} & \textbf{Назначение} \\
\midrule
Logistic Regression &
линейная supervised & нижняя граница линейного классификатора \\
Random Forest~\cite{breiman2001rf} &
ансамбль supervised & сильный baseline на табличных данных \\
XGBoost~\cite{chen2016xgboost} &
бустинг supervised & сильный baseline на табличных данных \\
MLP &
DL supervised & DL-baseline без учёта структуры \\
1D-CNN &
DL supervised & DL-baseline с локальной структурой \\
LSTM~\cite{hochreiter1997lstm} &
DL supervised & DL-baseline с временной структурой \\
GRU~\cite{cho2014gru} &
DL supervised & облегчённая RNN \\
BiLSTM &
DL supervised & альтернативная сильная DL-архитектура \\
TCN~\cite{bai2018tcn} &
DL supervised & альтернатива RNN \\
Transformer-encoder~\cite{vaswani2017attention} &
DL supervised & attention-based альтернатива \\
Autoencoder~\cite{mirsky2018kitsune} &
DL semi-supervised & контрольный AE \\
VAE~\cite{kingma2014vae, lopez2017cvae} &
DL semi-supervised & контрольный VAE \\
Isolation Forest~\cite{liu2008isolation} &
unsupervised & классический unsupervised \\
One-Class SVM~\cite{scholkopf2001ocsvm} &
unsupervised / one-class & классический one-class \\
\bottomrule
\end{tabularx}
\end{table}

\subsection{Сводная диаграмма экспериментального стенда}

\begin{figure}[H]
\centering
\begin{minipage}{0.94\linewidth}
\small
\begin{verbatim}
+----------------------------+        +----------------------------+
|  UNSW-NB15  (raw CSV)      |        |  CICIDS2017 (raw CSV)      |
+--------------+-------------+        +--------------+-------------+
               | preprocess + windowing              | preprocess + map
               v                                     v
+----------------------------+        +----------------------------+
|  Train / Val / Test         |       |  CICIDS2017 test (cross)   |
+--------------+--------------+       +--------------+--------------+
               | E1, E2, E3                          | E4'
               v                                     v
+--------------------------------------------------------------+
|              Models pool (Tab. 3.4): CNN-LSTM, ...           |
+----+--------------------+-----------------+------------+-----+
     | E1                 | E3              | E5        | E6
     v                    v                 v           v
+----------+        +----------+      +-----------+   +-----------+
| Compare  |        | Calibr.  |      | Drift inj |   | Lat./Thr. |
+----------+        +----------+      +-----------+   +-----------+
     ^                   ^                  ^             ^
     | E4                                                 |
     |                                                    |
+----+-----------------------------+                      |
|  AI Attacker (templates+FSM+drift) -------- stream ----+
+----------------------------------+
\end{verbatim}
\end{minipage}
\caption{Концептуальная схема экспериментального стенда}
\label{fig:experiment-stand}
\end{figure}

\subsection{Критерии успеха Этапа~4}

Этап~4 считается выполненным, если выполнены следующие условия.

\begin{enumerate}
    \item Полностью проведены эксперименты E1--E6, оформлены таблицы и
    графики метрик, графики обучения, confusion matrix.
    \item Гипотезы $H_0^{\text{main}}$, $H_0^{\text{aux}}$ и
    $H_0^{\text{transfer}}$ либо подтверждены при выполнении соответствующих
    NFR, либо обоснованно опровергнуты с указанием конкретных условий
    невыполнения.
    \item Сформирован итоговый отчёт с разделом ``Ограничения и направления
    дальнейшего развития'', явно фиксирующим невыполненные пункты и причины.
\end{enumerate}

\subsection{Промежуточные выводы по разделу}

\begin{enumerate}
    \item Сформирована библиотека из семи параметризованных шаблонов
    атакующих сценариев, согласованных с таксономией UNSW-NB15, с явной
    спецификацией параметров и привязкой каждого шаблона к классу датасета.
    \item Описаны конечно-автоматная политика агента, протокол валидации
    генератора (маргинальное соответствие нормы по KS и MMD;
    семантическая согласованность атак по средним и медианам).
    \item Сформирован протокол шести экспериментов (E1--E6 + E4') с
    явными целями, метриками и критериями принятия. Для всех экспериментов
    задана статистическая методология (множественные seed, bootstrap CI,
    парный критерий Уилкоксона).
    \item Зафиксирован пул из 14 моделей-кандидатов и baselines с
    указанием их роли в эксперименте E1; все модели обучаются и
    оцениваются в равных условиях, что устраняет ограничения, выявленные
    в опубликованной NIDS-литературе.
    \item Приведена диаграмма экспериментального стенда, отражающая
    взаимодействие подсистем: данные (UNSW-NB15, CICIDS2017),
    модели, активный тестовый источник (ИИ-злоумышленник) и блоки
    калибровки, drift-инжекции и измерения производительности.
\end{enumerate}


% ----------- Заключение -----------
\section*{ЗАКЛЮЧЕНИЕ}
\addcontentsline{toc}{section}{ЗАКЛЮЧЕНИЕ}

В результате выполнения Этапа~2 получены следующие основные результаты.

\begin{enumerate}
    \item Систематизированы опубликованные NIDS-результаты применения
    восьми типов архитектур нейронных сетей с явным указанием
    первоисточников, использованных датасетов и постановок задачи.
    Сводное сопоставление подтверждает выбор CNN-LSTM в качестве
    основной модели работы; альтернативами в собственном эксперименте
    останутся BiLSTM и Transformer-encoder, контрольными
    semi-supervised baseline --- AE и VAE.
    \item Зафиксированы методологические ограничения сравнений из
    опубликованной литературы (несовместимость датасетов, отсутствие
    контроля дрейфа, доминирование Accuracy при дисбалансе классов),
    что определяет необходимость собственного эксперимента в равных
    условиях.
    \item Полностью описана теоретическая модель данных и pipeline:
    схема UNSW-NB15, подмножество признаков, общих с CICIDS2017,
    стадии очистки, кодирования (one-hot, frequency, target, embedding),
    масштабирования (log, robust), формирования окон (\(W=32\),
    \(S=W/4\)), методы управления дисбалансом (re-weighting, focal loss,
    SMOTE), блок drift-monitoring (PSI, KL, MMD).
    \item Сформирована библиотека из семи параметризованных шаблонов
    атак, согласованных с таксономией UNSW-NB15, и описана
    конечно-автоматная политика ИИ-злоумышленника. Зафиксирован
    протокол валидации генератора (маргинальное соответствие нормы по
    тесту Колмогорова--Смирнова и MMD; семантическая согласованность
    атак).
    \item Сформирован протокол шести экспериментов (E1--E6, плюс E4'
    cross-evaluation на CICIDS2017) с явными целями, метриками и
    количественными критериями принятия гипотез. Для каждого эксперимента
    задана статистическая методология (множественные seed, bootstrap CI,
    парный критерий Уилкоксона).
\end{enumerate}

Полученные результаты используются как входные данные Этапа~3 ВКР:
программная реализация прототипа, обучение моделей, реализация
ИИ-злоумышленника и подготовка стенда экспериментов.

% ----------- Список использованных источников -----------
\printbibliography[title={СПИСОК ИСПОЛЬЗОВАННЫХ ИСТОЧНИКОВ}]

\label{LastPage}

\end{document}
